\documentclass[a4paper]{article}

\usepackage{graphicx}
\usepackage{amsmath,amssymb}
\usepackage{minted}
\usepackage{multirow}
\usepackage{float}
\usepackage{algorithm}
\usepackage{algpseudocode}
\usepackage{todonotes}
\usepackage{forest}
\usepackage{xstring}
\usepackage{xcolor}
\usepackage[most]{tcolorbox}
\usepackage{xparse}
\usepackage{natbib}

\usetikzlibrary{graphs,graphdrawing}
\usegdlibrary{layered}

% for external graphviz
\input{/home/donkarlo/Dropbox/repo/nd_language_ai_project/src/nd_language_ai/formal/computer/programming/tex/latex/code_clips/external_graphviz.tex}

% for tags
\input{/home/donkarlo/Dropbox/repo/nd_language_ai_project/src/nd_language_ai/formal/computer/programming/tex/latex/part/control_sequence/kind/semtag/semtag.tex}

% for links
\input{/home/donkarlo/Dropbox/repo/nd_language_ai_project/src/nd_language_ai/formal/computer/programming/tex/latex/code_clips/seehere.tex}

% for nested lists and sections
\input{/home/donkarlo/Dropbox/repo/nd_language_ai_project/src/nd_language_ai/formal/computer/programming/tex/latex/code_clips/deep_structure.tex}

\title{Startup Regression after Navigator Context-Menu Consolidation}
\date{2026-09-06}
\author{Mohammad Rahmani}

\begin{document}
   \maketitle

   \section{Failure}
   Mind Mirror exited during import because the legacy
   \texttt{Fix11DesktopController} tried to capture
   \texttt{WorkspaceInteractionController.\_show\_navigator\_menu}, while the
   newer Navigator consolidation intentionally removed that secondary menu
   implementation.

   \section{Correction}
   The legacy controller no longer patches the Workspace interaction menu.
   Hierarchical search, exact image drag/drop, and generic iPad image editing
   remain installed there.  The canonical
   \texttt{NavigatorMultiPathContextController} now exposes
   \texttt{Edit image in iPad} for all supported raster formats as well as
   managed \texttt{.ndgraphic} files, while retaining Ctrl/Shift
   multi-selection and Copy/Cut/Paste.

   \section{Regression coverage}
   A regression test verifies that startup code contains no dependency on the
   removed Workspace menu method, and that the canonical multi-path menu owns
   both file clipboard actions and iPad image editing.

   \bibliographystyle{plainnat}
   \bibliography{/home/donkarlo/Dropbox/repo/nd_shared_research/bibliography/bibliography.bib}
\end{document}
